\chapter{Introdução}
\label{sec:introducao}

\section{Contextualização e Motivação}
O monitoramento de condições ambientais em espaços internos desempenha um papel importante na caracterização e no acompanhamento das condições de funcionamento desses ambientes. Entre as variáveis frequentemente monitoradas, a temperatura se destaca por sua relação com o conforto térmico, com o funcionamento de equipamentos e com as condições de operação de determinados espaços. Nesse contexto, a utilização de dispositivos conectados à Internet das Coisas (Internet of Things — IoT) permite realizar medições periódicas e armazenar grandes volumes de dados para posterior análise.

Sistemas de monitoramento baseados em sensores distribuídos possibilitam acompanhar a evolução temporal da temperatura em diferentes posições de um ambiente. Entretanto, as medições obtidas por esses dispositivos representam apenas as condições observadas em seus respectivos pontos de instalação. A distribuição espacial da temperatura pode não ser uniforme, especialmente em ambientes que apresentam equipamentos eletrônicos, sistemas de climatização, divisórias e diferentes condições de circulação de ar. Dessa forma, a temperatura registrada em um determinado ponto não necessariamente representa a temperatura de todo o espaço monitorado.

Uma alternativa para ampliar o aproveitamento das informações disponíveis é o emprego de técnicas de sensoriamento virtual, como discutido por \citeonline{brunello2021virtual}. Diferentemente de um sensor físico, que realiza diretamente a medição de uma grandeza, um sensor virtual utiliza informações provenientes de outras fontes para estimar o valor dessa grandeza. No monitoramento térmico, essa abordagem permite investigar a possibilidade de estimar a temperatura em uma determinada posição a partir das medições realizadas por sensores instalados em outras posições.

A viabilidade dessa abordagem depende, entre outros fatores, das relações existentes entre as séries temporais dos sensores, da distribuição espacial dos dispositivos e do método utilizado para realizar as estimativas. Sensores próximos podem apresentar comportamentos semelhantes, mas essa relação não é necessariamente determinada apenas pela distância física. Características específicas do ambiente também podem influenciar as medições e, consequentemente, a precisão das estimativas.

Nesse cenário, métodos estatísticos e de aprendizado de máquina podem ser empregados para identificar relações entre as temperaturas registradas por diferentes dispositivos. Modelos como regressão linear e Random Forest permitem construir estimativas a partir de observações históricas, oferecendo alternativas a estratégias mais simples, como o cálculo da média das temperaturas dos sensores disponíveis.

O presente trabalho investiga a reconstrução virtual de medições de temperatura utilizando dados históricos de sensores IoT instalados em ambientes internos de laboratório. O estudo considera uma infraestrutura de monitoramento já existente, sem alterações no posicionamento físico dos dispositivos. Dessa forma, a investigação concentra-se no aproveitamento dos dados disponíveis e na avaliação de diferentes estratégias de estimativa.

Para possibilitar a avaliação quantitativa dos métodos, adota-se uma abordagem na qual as medições de um sensor físico são utilizadas como referência, enquanto as medições dos demais sensores servem como entradas para os modelos. Assim, é possível comparar as temperaturas estimadas com os valores efetivamente registrados e quantificar os erros de reconstrução.

Além da comparação entre diferentes modelos, o trabalho busca investigar como a quantidade de sensores utilizados e os critérios empregados para sua seleção influenciam a qualidade das estimativas. Para isso, são consideradas tanto as relações estatísticas entre as séries temporais quanto as informações aproximadas sobre a localização física dos dispositivos.

\section{Delimitação e Evolução do Escopo}
\label{sec:evolucao_escopo}

A proposta inicial, intitulada \textit{Enriquecimento de Dados de Sensores
IoT para Aprendizado de Máquina focado em Conforto Térmico no LASDPC},
previa estruturar os históricos e associá-los à localização física,
aos estados da climatização e a outras informações para apoiar modelos
de estimativa térmica e inferência de conforto. O presente documento
explicita o recorte efetivamente executado: a reconstrução de temperaturas
registradas por sensores físicos, com referência observada para avaliação.

A preparação dos dados resultou em matrizes temporais sincronizadas.
O enriquecimento foi operacionalizado pela associação dos identificadores
ao cadastro e às coordenadas, e pela integração da umidade às entradas
dos modelos. A regressão linear e a Random Forest constituem as abordagens
de aprendizado supervisionado avaliadas. As coordenadas apoiam a seleção
por proximidade e a representação espacial; não foram adicionadas como
variáveis explicativas dos modelos individuais de cada alvo.

O conforto térmico permanece como motivação e possibilidade de aplicação.
Não foram coletadas todas as variáveis necessárias a uma avaliação de
conforto, nem respostas dos ocupantes que permitissem validá-la. Também
não foram integradas informações meteorológicas externas ou estados
confirmados dos aparelhos de climatização. Os registros de energia foram
limitados a uma auditoria, devido à curta disponibilidade e à unidade
não confirmada do campo consultado.

Assim, a contribuição apresentada está na estruturação, no enriquecimento
e na avaliação dos dados para sensoriamento virtual. O mapa permite
visualizar diferenças entre temperaturas médias em posições conhecidas,
mas não representa a validação de um gradiente térmico contínuo ou de
condições de conforto. A inferência de conforto e o controle automatizado,
previstos como aplicações na proposta inicial, dependem de uma etapa
adicional de instrumentação e definição do protocolo de avaliação.

\section{Objetivos}
No recorte delimitado, o objetivo geral é avaliar a reconstrução virtual
de temperaturas a partir de dados históricos de sensores IoT,
estruturados e enriquecidos com cadastro espacial e umidade,
utilizando medições físicas como referência experimental.

Como objetivos específicos, estabelecem-se:
\begin{itemize}
    \item caracterizar os dados históricos e a disponibilidade das medições;
    \item organizar os sensores por identificadores cadastrais e posições relativas;
    \item preparar e sincronizar as séries em uma grade temporal comum;
    \item comparar média aritmética, regressão linear e Random Forest;
    \item investigar o efeito da quantidade de sensores de entrada;
    \item comparar seleção por correlação e proximidade no subconjunto com coordenadas;
    \item avaliar a umidade dos sensores de entrada como informação complementar
    à reconstrução de temperatura;
    \item interpretar os erros por sensor e representar espacialmente as temperaturas médias.
\end{itemize}

\section{Organização do Trabalho}
Este trabalho está organizado em quatro capítulos. O
Capítulo~\ref{sec:revisao} apresenta os conceitos de sensoriamento virtual,
séries temporais e métodos de estimativa, além dos trabalhos relacionados.
O Capítulo~\ref{sec:desenvolvimento} descreve a infraestrutura, os dados,
o pré-processamento e o protocolo de avaliação, e apresenta os resultados
e suas limitações. O Capítulo~\ref{sec:conclusao} reúne as conclusões e as
possibilidades de continuidade da pesquisa.
